\section{Models}
\label{sec:models}

\paragraph{Simple signals.}
Four signals need no fitting: 12--1 momentum~\cite{jegadeesh1993returns}, minus the past five-day return (reversal)~\cite{jegadeesh1990evidence}, the standardized earnings surprise (SUE)~\cite{livnat2006comparing}, and the three-month revision of the consensus forecast scaled by price~\cite{chan1996momentum}.

\begin{table}[t]
\caption{Features. Each value at date $t$ only uses data available at the close of $t$; stock-level variables are replaced by their centred percentile rank within each date over the training universe.}
\label{tab:features}
\footnotesize
\begin{tabular}{@{}p{0.25\linewidth}p{0.71\linewidth}@{}}
\toprule
Group & Variables \\
\midrule
Past returns (11) & log returns over 1, 5, 10, 21, 63, 126 and 252 days; 12--1 momentum; gaps to the 20-, 50- and 200-day moving averages of the log price \\
Risk (8) & drawdown from the 252-day high; volatility over 5, 21 and 63 days; largest daily return over 21 days; 63-day skewness; 252-day market beta; 63-day idiosyncratic volatility \\
Liquidity and size (6) & log 21-day dollar volume; 5- to 63-day volume ratio; 21-day share turnover; 21-day Amihud illiquidity; log market value; log listing age \\
Market (3) & market return over 5 and 21 days; 21-day market volatility (not ranked) \\
Accounting (13) & return on assets and on equity; gross profitability; leverage; cash to assets; current ratio; asset, sales and earnings growth; book-to-market; earnings, sales and EBITDA yields \\
Analysts (7) & consensus EPS to price; its 1- and 3-month revisions; balance of upward and downward revisions; dispersion; log number of analysts; standardized earnings surprise \\
\bottomrule
\end{tabular}
\end{table}

\paragraph{Tabular models.}
The market features are 25 stock-level variables (returns from 1 to 252 days, gaps to moving averages, volatility, beta, idiosyncratic volatility, liquidity and size) and 3 market-level ones.
We add 13 accounting variables (Acc) and 7 analyst variables (Ana), listed in \cref{tab:features}.
Stock-level variables are replaced by their centred percentile rank within each date.
LightGBM~\cite{ke2017lightgbm} (learning rate 0.02, 50\% feature and row subsampling, L2 penalty 10) has a grid of 7 or 31 leaves, 500 or 5,000 minimum samples per leaf and 50 to 400 rounds, and one setting is \emph{selected} each year on validation.
Because the selection changed in almost every year, we also test a \emph{grid average}: the four tree settings are fit on all past data, and the within-date ranks of their predictions after 100, 200 and 400 rounds are averaged (12 models).

\paragraph{TimesFM~3 zero-shot.}
TimesFM~3~\cite{timesfm3} is a decoder-only transformer with 330.7M parameters that reads 32-step input patches and writes 64-step output patches with quantiles from 0.1 to 0.9.
In univariate mode, each forecast reads up to the last 512 daily returns of a stock, with missing days set to zero.
The score is the forecast cumulative return over $t{+}2$ to $t{+}6$; we also test the log price index as input.
The forecast width and asymmetry, added with the score as ranked LightGBM features, are computed over the training universe.
The width tracks the realized volatility of the next five days (rank correlation 0.60), which confirms that the model is fed correctly.

\paragraph{Fine-tuning.}
The package documents only inference and its decoding blocks gradients; calling the undecorated function reproduces the inference outputs exactly.
We train the last four transformer layers and the output head (66.3M parameters).
Each step draws one training date and 64 stocks of the training universe with a full 512-day context.
The score $s_i$ sums the median forecasts over $t{+}2$ to $t{+}6$, and the loss is minus its correlation with the realized sum of returns $y_i$:
\begin{equation}
\mathcal{L} = -\frac{\sum_i (s_i-\bar s)(y_i-\bar y)}{\lVert s-\bar s\rVert\,\lVert y-\bar y\rVert},
\label{eq:loss}
\end{equation}
which rewards the ranking within a date and ignores the market level.
We use AdamW (learning rate $3\times10^{-5}$, weight decay 0.01) with gradient clipping at 1, and skip the steps whose gradient is not finite (about 3\%, from a square root in the model's input normalization).
These settings, with 1,000 steps, were fixed on data before 2005, so no development year was used to tune them.
The model is retrained every three years on all windows labelled before each block, with five seeds whose within-date ranks are averaged.
One follow-up, declared in advance after the development results, changes only the number of steps to 10,000, which takes about 48 minutes per run on one RTX A4500 GPU.

\paragraph{Ensembles.}
Ensembles average the within-date percentile ranks of their members, with equal weights and nothing fitted, and every ensemble was listed before its results were computed.
The \textbf{primary model} gives equal weights to the grid-averaged LightGBM with market, accounting and analyst features (averaged over three seeds) and to the four simple signals.
